%%%%%%%%%%%%%%%%%%%%%%%%%%%%%%%%%%%%%%%%%
% Short Sectioned Assignment LaTeX Template Version 1.0 (5/5/12)
% This template has been downloaded from: http://www.LaTeXTemplates.com
% Original author:  Frits Wenneker (http://www.howtotex.com)
% License: CC BY-NC-SA 3.0 (http://creativecommons.org/licenses/by-nc-sa/3.0/)
%%%%%%%%%%%%%%%%%%%%%%%%%%%%%%%%%%%%%%%%%

% \documentclass[paper=a4, fontsize=11pt]{scrartcl} % A4 paper and 11pt font size
\documentclass[11pt, a4paper]{book}
\usepackage[T1]{fontenc} % Use 8-bit encoding that has 256 glyphs
\usepackage[utf8]{inputenc}
\usepackage{fourier} % Use the Adobe Utopia font for the document - comment this line to return to the LaTeX default
\usepackage{listings} % para insertar código con formato similar al editor
\usepackage[spanish, es-tabla]{babel} % Selecciona el español para palabras introducidas automáticamente, p.ej. "septiembre" en la fecha y especifica que se use la palabra Tabla en vez de Cuadro
\usepackage{url} % ,href} %para incluir URLs e hipervínculos dentro del texto (aunque hay que instalar href)
\usepackage{graphics,graphicx, float} %para incluir imágenes y colocarlas
\usepackage[gen]{eurosym} %para incluir el símbolo del euro
\usepackage{cite} %para incluir citas del archivo <nombre>.bib
\usepackage{enumerate}
\usepackage{hyperref}
\usepackage{graphicx}
\usepackage{tabularx}
\usepackage{booktabs}

\usepackage[table,xcdraw]{xcolor}
\hypersetup{
	colorlinks=true,	% false: boxed links; true: colored links
	linkcolor=black,	% color of internal links
	urlcolor=cyan		% color of external links
}
\renewcommand{\familydefault}{\sfdefault}
\usepackage{fancyhdr} % Custom headers and footers
\pagestyle{fancyplain} % Makes all pages in the document conform to the custom headers and footers
\fancyhead[L]{} % Empty left header
\fancyhead[C]{} % Empty center header
\fancyhead[R]{Daniel González Serrano} % My name
\fancyfoot[L]{} % Empty left footer
\fancyfoot[C]{} % Empty center footer
\fancyfoot[R]{\thepage} % Page numbering for right footer
%\renewcommand{\headrulewidth}{0pt} % Remove header underlines
\renewcommand{\footrulewidth}{0pt} % Remove footer underlines
\setlength{\headheight}{13.6pt} % Customize the height of the header

\usepackage{titlesec, blindtext, color}
\definecolor{gray75}{gray}{0.75}
\newcommand{\hsp}{\hspace{20pt}}
\titleformat{\chapter}[hang]{\Huge\bfseries}{\thechapter\hsp\textcolor{gray75}{|}\hsp}{0pt}{\Huge\bfseries}
\setcounter{secnumdepth}{4}
\usepackage[Lenny]{fncychap}


\begin{document}

	% Plantilla portada UGR
	\begin{titlepage}
\newlength{\centeroffset}
\setlength{\centeroffset}{-0.5\oddsidemargin}
\addtolength{\centeroffset}{0.5\evensidemargin}
\thispagestyle{empty}

\noindent\hspace*{\centeroffset}\begin{minipage}{\textwidth}

\centering
\includegraphics[width=0.9\textwidth]{logos/logo_ugr.jpg}\\[1.4cm]

\textsc{ \Large TRABAJO FIN DE GRADO\\[0.2 cm]}
\textsc{ GRADO EN INGENIERÍA INFORMÁTICA}\\[1 cm]

{\Huge\bfseries Memes para todos \\}
\noindent\rule[-1ex]{\textwidth}{3 pt}\\[3.5 ex]
{\large\bfseries Almacenamiento Online }
\end{minipage}

\vspace{2.5cm}
\noindent\hspace*{\centeroffset}
\begin{minipage}{\textwidth}
\centering

\textbf{Autor}\\ {Daniel González Serrano}\\[2.5 ex]
\textbf{Director}\\ {Juan Julián Merelo Guervós}\\[2 cm]
\includegraphics[width=0.3\textwidth]{logos/etsiit_logo.png}\\[0.1 cm]
\textsc{Escuela Técnica Superior de Ingenierías Informática y de Telecomunicación}\\
\textsc{---}\\
Granada, \today
% Mostrar nombre de mes y año de entrega del documento

\end{minipage}
\end{titlepage}


	% Plantilla prefacio UGR
	\thispagestyle{empty}

\begin{center}
{\large\bfseries Memes para todos \\ Almacenamiento Online }\\
\end{center}
\begin{center}
Daniel González Serrano\\
\end{center}

%\vspace{0.7cm}

\vspace{0.5cm}
\noindent\textbf{Palabras clave}: \textit{software libre}
\vspace{0.7cm}

\noindent\textbf{Resumen}\\
	

\cleardoublepage

\begin{center}
	{\large\bfseries Igual, pero en inglés}\\
\end{center}
\begin{center}
	Daniel González Serrano\\
\end{center}
\vspace{0.5cm}
\noindent\textbf{Keywords}: \textit{open source}, \textit{floss}
\vspace{0.7cm}

\noindent\textbf{Abstract}\\


\cleardoublepage

\thispagestyle{empty}

\noindent\rule[-1ex]{\textwidth}{2pt}\\[4.5ex]

D. \textbf{Juan Julián Merelo Guervós}, Profesor(a) del ...

\vspace{0.5cm}

\textbf{Informo:}

\vspace{0.5cm}

Que el presente trabajo, titulado \textit{\textbf{Chief}},
ha sido realizado bajo mi supervisión por \textbf{Estudiante}, y autorizo la defensa de dicho trabajo ante el tribunal
que corresponda.

\vspace{0.5cm}

Y para que conste, expiden y firman el presente informe en Granada a \today

\vspace{1cm}

\textbf{El/la director(a)/es: }

\vspace{5cm}

\noindent \textbf{(nombre completo tutor/a/es)}

\chapter*{Agradecimientos}






	% Índice de contenidos
	\newpage
	\tableofcontents

	% Índice de imágenes y tablas
	\newpage
	\listoffigures

	% Si hay suficientes se incluirá dicho índice
	\listoftables 
	\newpage

	% Introducción 
	\chapter{Introducción}

Este proyecto es software libre, y está liberado con la licencia \cite{gplv3}.

	% Estado del arte
	\chapter{Estado del Arte}

En el ámbito del almacenamiento online, se encuentran diversas soluciones
que abordan la gestión de contenido digital en especial, imágenes.

\section{Crítica al estado del arte}

Se han identificado aplicaciones como Dropbox \cite{dropbox}, Google Fotos \cite{googlefotos}, OneDrive \cite{onedrive}, iCloud \cite{icloud}, Amazon Photos \cite{amazonphotos} con capacidad de colaboración y almacenamiento.
Además, se han encontrado alternativas más orientadas a la compartición de imágenes como Flickr \cite{flickr}, Imgur \cite{imgur}, 500px \cite{500px}, entre otros que también permiten subir imágenes a la nube.

A continuación, se presenta una tabla comparativa de las aplicaciones más destacadas de las mencionadas anteriormente. En las columnas se encuentran las características que se han considerado más relevantes para el desarrollo de la aplicación:

\begin{itemize}
    \item \textbf{Almacenamiento gratuito:} Capa gratuita de almacenamiento.
    \item \textbf{Organización:} Característica diferenciadora para organizar.
    \item \textbf{Multiplataforma:} Capacidad de funcionar desde diferentes dispositivos.
    \item \textbf{Colaborativo:} Capacidad de compartir contenido con otros usuarios.
    \item \textbf{Seguridad:} Mecanismos de seguridad y privacidad como encriptación de contenido o carpetas seguras.
    \item \textbf{Edición:} Capacidad de editar imágenes.
    \item \textbf{Control de versiones:} Capacidad de mantener versiones de imágenes.
    \item \textbf{Precio:} Costo de la aplicación.
\end{itemize}

% Google Fotos
\begin{tabular}{| c | c | c | c |}
    \hline
    \cellcolor[HTML]{4285F3} & Google Fotos\\ 
    \hline
    Almacenamiento gratuito & 15 GB \\
    \hline
    Organización & Automática e inteligente \\
    \hline
    Multiplataforma & Sí \\
    \hline
    Colaborativo & Sí \\
    \hline
    Seguridad & Cifrado, carpeta privada \\
    \hline
    Edición & Sí \\
    \hline
    Control de versiones & No \\
    \hline
    Precio & 100 GB - 1,99 €/mes \\
    \hline
\end{tabular}

\vspace{1cm} % Espacio entre tablas

% OneDrive
\begin{tabular}{| c | c | c | c |}
    \hline
        \cellcolor[HTML]{24a0e8} & OneDrive\\ 
    \hline
    Almacenamiento gratuito & 5 GB \\
    \hline
    Organización & Álbumes automáticos \\
    \hline
    Multiplataforma & Sí \\
    \hline
    Colaborativo & Sí \\
    \hline
    Seguridad & Cifrado, bóveda personal \\
    \hline
    Edición & No \\
    \hline
    Control de versiones & Sí \\
    \hline
    Precio & 100 GB (Microsoft 365) - 20 €/año \\
    \hline
\end{tabular}

\vspace{1cm} % Espacio entre tablas

% Amazon Photos
\begin{tabular}{| c | c | c | c |}
    \hline
        \cellcolor[HTML]{00464f}& Amazon Photos\\ 
    \hline
    Almacenamiento gratuito & 5 GB \\
    \hline
    Organización & Mediante IA \\
    \hline
    Multiplataforma & Sí \\
    \hline
    Colaborativo & Sí \\
    \hline
    Seguridad & Cifrado \\
    \hline
    Edición & No \\
    \hline
    Control de versiones & No \\
    \hline
    Precio & Ilimitado - Prime 4,99 €/mes \\
    \hline
\end{tabular}

\vspace{1cm} % Espacio entre tablas

% iCloud
\begin{tabular}{| c | c | c | c |}
    \hline
    \cellcolor[HTML]{000000} & iCloud (Se necesita un Apple ID)\\ 
    \hline
    Almacenamiento gratuito & 5 GB \\
    \hline
    Organización & Reconocimiento facial \\
    \hline
    Multiplataforma & Sí (Solo Apple) \\
    \hline
    Colaborativo & Sí (Solo Apple) \\
    \hline
    Seguridad & Cifrado, carpeta segura \\
    \hline
    Edición & No \\
    \hline
    Control de versiones & No \\
    \hline
    Precio & 50 GB - 0.99 €/mes \\
    \hline
\end{tabular}

\vspace{1cm} % Espacio entre tablas

% MEGA
\begin{tabular}{| c | c | c | c |}
    \hline
        \cellcolor[HTML]{dd1405} & MEGA\\ 
    \hline
    Almacenamiento gratuito & 20 GB \\
    \hline
    Organización & Básica \\
    \hline
    Multiplataforma & Sí \\
    \hline
    Colaborativo & Sí \\
    \hline
    Seguridad & Cifrado + contraseña \\
    \hline
    Edición & No \\
    \hline
    Control de versiones & Sí \\
    \hline
    Precio & 2 TB - 9,99 €/mes \\
    \hline
\end{tabular}

\vspace{1cm} % Espacio entre tablas

% Dropbox
\begin{tabular}{| c | c | c | c |}
    \hline
        \cellcolor[HTML]{0061fe} & Dropbox\\ 
    \hline
    Almacenamiento gratuito & 2 GB \\
    \hline
    Organización & Mediante etiquetas (solo web) \\
    \hline
    Multiplataforma & Sí \\
    \hline
    Colaborativo & Sí \\
    \hline
    Seguridad & Cifrado + contraseña \\
    \hline
    Edición & No \\
    \hline
    Control de versiones & Sí \\
    \hline
    Precio & 2 TB - 11,99 €/mes \\
    \hline
\end{tabular}

\section{Propuesta}
Después del estudio del panorama actual relacionado con el almacenamiento online de contenido digital, se ha identificado el problema de la inexistencia de un sistema centralizado que permita a los usuarios almacenar y compartir imágenes de forma segura y sencilla.

El sistema propuesto se centrará en la gestión de imágenes, ofreciendo un almacenamiento en la nube, con la posibilidad de compartir imágenes con otros usuarios y la capacidad de editarlas, todo, desde cualquier dispositivo. Además, se incluirá un sistema de organización automática e inteligente y la posibilidad de mantener versiones de las imágenes.

	% Planificación
	\chapter{Planificación}

\section{Temporización}

\section{Seguimiento de la Temporización}
	

	% Metodología
	\chapter{Metodología}

\section{Enfoque Ágil}

\section{Design Thinking}

\section{Herramientas y Teconologías}


	% Implementación
	\chapter{Implementación}
 
\section{Domain Driven Design}

\section{Diseño e Implementación Iterativa}

\section{Validación de Progreso}

\section{Tecnologías}

	% Tests y despliegue
	\chapter{Tests y despliegue}

\section{Pruebas unitarias}

\section{Pruebas de integración}

\section{Despliegue}

	% Costes
	\chapter{Costes}

\section{Infraestructura}

\section{Proyecto}



	% Trabajos futuros
	\chapter{Análisis de Prestaciones}

\section{Pruebas de Carga}

\section{Resultados y Comportamiento}




	% Conclusiones
	\chapter{Conclusiones y trabajos futuros}

Finalmente, en este capítulo se presentarán las conclusiones obtenidas tras la realización del proyecto y se propondrán posibles trabajos futuros que se podrían llevar a cabo para mejorar la aplicación

El desarrollo de este proyecto ha permitido alcanzar los objetivos planteados inicialmente. Como se estableció en~\ref{sec:objetivos_iniciales}, se diseñó una \textbf{solución económica y accesible}, \textbf{adaptable a una amplia gama de clientes} y \textbf{compatible con dispositivos de bajas prestaciones}. Esto se logró utilizando tecnologías eficientes como Next.js y Fabric.js que, además de reducir costos, optimizan la accesibilidad y la experiencia del usuario final en diversas plataformas. Asimismo, el uso de una \textbf{licencia abierta} como es la GPLv3~\cite{gplv3} permite que la aplicación sea utilizada sin restricciones, promoviendo su uso en entornos personales, educativos y empresariales.

Por otra parte, las historias de usuario que se marcaron al principio del proyecto se han cumplido en su totalidad, aunque una de ellas no se ha podido implementar al completo. Así pues, se va a proceder a enumerar las historias de usuario y a describir su estado actual:

\begin{itemize}
  \item[\bien] \textbf{HU01 Como Cristina Contreras Márquez, estudiante de Ingeniería Informática, como hago memes mientras voy haciendo las presentaciones quiero rapidez y sencillez por lo que no quiero tener que hacer el meme desde cero:} Esta historia de usuario se ha satisfecho completamente. Cristina Contreras Márquez puede acceder a algún catálogo hecho con memes realizados por otro usuario y modificarlo a su gusto.
  \item[\bien] \textbf{HU02 Como Cristina Contreras Márquez, estudiante de Ingeniería Informática, quiero crear memes que puedan ser modificados o reutilizados por otros usuarios.} De igual manera, Cristina Contreras Márquez puede elaborar memes y subirlos a la plataforma para que otros usuarios puedan modificarlos.
  \item[\bien] \textbf{HU03 Como Ted Johnson González, experto conferenciante, quiero acceso rápido y sencillo a memes relevantes y atractivos para incluir en mis presentaciones sin necesidad de confeccionarlos yo mismo.} En este caso, Ted Johnson González puede acceder a catálogos de memes que otros usuarios han elaborado y puede descargarlos para incluirlos en sus presentaciones.
  \item[\regular] \textbf{HU04 Como departamento de marketing, queremos colaborar en un meme varios miembros del equipo.} Por último, tenemos al departamento de marketing cuya historia de usuario no se ha cumplido al completo dado que no se ha implementado la funcionalidad de permisos, roles y colaboración en la aplicación, pero esto no impide que varios miembros del equipo puedan colaborar en un meme.
\end{itemize}

Por lo que acabamos de ver se puede afirmar que el proyecto ha sido un éxito en cuanto a la implementación de las historias de usuario marcadas al principio del proyecto. A continuación se presentarán las conclusiones obtenidas tras la realización del proyecto.

\section{Conclusiones}

Tras haber trabajado en este proyecto con últimas tecnologías como es Next.js 15 que salió este pasado 21 de octubre de 2024, he llegado a la conclusión que este marco de desarrollo web de React otorga al programador una experiencia de desarrollo muy agradable, ágil y eficiente. Además, la casi cero configuración que requiere el despliegue de la aplicación en Vercel cierra el flujo de trabajo de cualquier proyecto software de manera muy satisfactoria.

En cuanto a la otra tecnología ‘troncal’ utilizada como es Fabric.js 5.0, he de decir que me ha sorprendido gratamente la facilidad con la que se pueden crear lienzos interactivos y la cantidad de funcionalidades que ofrece. Al principio, cuesta bastante acostumbrarse a esta nueva forma de trabajar con elementos gráficos, pero es una librería muy potente y versátil.

A título personal siempre he querido experimentar el desarrollo de una aplicación web no solo con las últimas tecnologías sino también con enfoques ágiles y debo decir que ha sido una experiencia muy enriquecedora y que desde luego voy a poder aplicar tanto en el ámbito profesional como en futuros proyectos personales.

\section{Trabajos futuros}

Como se ha mencionado en la introducción de este capítulo la última historia de usuario (HU04) no se ha implementado al completo. Por lo que se propone como trabajo futuro la implementación de la funcionalidad de permisos, roles y colaboración en la aplicación. De esta manera, varios miembros del equipo podrán colaborar en un meme y se podrán establecer roles de usuario como administrador, editor, colaborador, etc.

Como se ha hecho durante toda la planificación de la aplicación se va a establecer un quinto hito o milestone contemplando todos esos trabajos futuros.

\subsection{Milestone 5: Autenticación, colaboración y roles}

El PMV u objetivo de este milestone es la implementación de los roles de usuario respecto a la modificación de memes, su respectiva lógica de colaboración y permisos así como el sistema de autenticación y perfiles.

Algunos de los objetivos que se pretenden alcanzar en este milestone son:

\begin{itemize}
  \item \textbf{Roles de usuario.}
  \item \textbf{Lógica de colaboración y permisos.}
  \item \textbf{Autenticación.}
  \item \textbf{Perfiles.}
\end{itemize}


	
	\newpage

	\bibliography{bibliografia}
	\bibliographystyle{plain}
	
\end{document}

